\chapter{Haag--Ruelle Theory and Mathematical Scattering}
\label{ch:volume-iv-haag-ruelle-mathematical-scattering}

\section{Scattering as a Theorem About Local Nets}

The earlier construction of the \(S\)-matrix used Haag--Ruelle limits before
LSZ and before perturbative scattering rules.  In local-net language the same
construction can be stated directly for a local net.
Let \(O\mapsto\mathcal R(O)\) be a Poincare-covariant vacuum net on
\(\Hilb\), with vacuum vector \(\Omega\), positive-energy translation
generators \(P^\mu\), and bounded local algebras.  A single massive particle
species of mass \(m>0\) is represented by an isolated spectral subspace
\[
  \Hilb_1=E(\Sigma_m^+)\Hilb,
  \qquad
  \Sigma_m^+=\{p:p^0=\sqrt{\vec p^{\,2}+m^2},\ p^2=-m^2\},
\]
where \(E\) is the joint spectral measure of \(P^\mu\).  Isolation means that
there is an open neighborhood \(N\) of \(\Sigma_m^+\) in momentum space such
that
\[
  \operatorname{Spec}(P)\cap N=\Sigma_m^+ .
\]

An operator \(B\) is almost local when there are local operators \(B_R\)
localized in double cones of radius \(R\) such that
\[
  \|B-B_R\|\le C_N R^{-N}
  \qquad\text{for every }N .
\]
Almost-local operators are obtained by smearing local observables in spacetime
with Schwartz functions whose Fourier transforms select the isolated shell.
They are the bounded-operator substitute for interpolating fields.

\section{Energy-Momentum Transfer and Regular Creators}

The scattering construction uses local observables only after their
energy-momentum transfer has been restricted.  For a bounded operator \(A\)
set
\[
  A(x)=U(x)AU(x)^{-1},
  \qquad x\in\R^D .
\]
For \(f\in\mathcal S(\R^D)\), define
\[
  A(f)=\int\dd^Dx\,f(x)A(x),
  \qquad
  \widehat f(q)=\int\dd^Dx\,f(x)\ee^{\ii q\cdot x}.
\]
With this Fourier convention, \(A(f)\) changes the joint spectral support of
the translation generators only by momenta in \(\operatorname{supp}\widehat f\):
\[
  E(\Delta_1)A(f)E(\Delta_2)=0
  \quad\text{if}\quad
  \Delta_1\cap\bigl(\Delta_2+\operatorname{supp}\widehat f\bigr)=\varnothing .
\]
Here \(\Delta_1,\Delta_2\subset\R^D\) are Borel sets and \(E\) is the joint
spectral measure of \(P^\mu\).  This is the spectral-transfer statement
behind the phrase ``selecting the one-particle shell.''

Choose a compact set \(K\subset\Sigma_m^+\) and a Schwartz function \(\chi\)
whose Fourier transform is supported in the isolating neighborhood \(N\) and
is equal to \(1\) on a smaller neighborhood of \(K\).  If \(A\) is local and
\[
  B=A(\chi),
\]
then \(B\) is almost local.  Since \(\Omega\) has momentum \(0\), the vector
\(B\Omega\) has spectral support contained in
\(\operatorname{supp}\widehat\chi\cap\operatorname{Spec}(P)\).  By the
isolation hypothesis this intersection lies on \(\Sigma_m^+\), and therefore
\[
  B\Omega=P_1B\Omega .
\]
The adjoint \(B^*\) has transfer contained in
\(-\operatorname{supp}\widehat\chi\).  If this reflected support is disjoint
from the forward spectrum except at the origin, and the support of
\(\widehat\chi\) avoids \(0\), then
\[
  B^*\Omega=0 .
\]
An almost-local operator satisfying these spectral properties and
\(B\Omega\ne0\) will be called a regular Haag--Ruelle creator for the chosen
mass shell.

\begin{figure}[H]
  \centering
  \resizebox{0.78\textwidth}{!}{%
  \begin{tikzpicture}[x=1cm,y=1cm,every node/.style={font=\scriptsize}]
    \draw[-{Latex[length=2mm]},line width=0.75pt] (-3.2,0)--(3.45,0)
      node[right] {\(\vec p\)};
    \draw[-{Latex[length=2mm]},line width=0.75pt] (0,-2.35)--(0,3.2)
      node[above] {\(p^0\)};
    \draw[gray!50,dashed,line width=0.8pt] (-2.9,2.9)--(0,0)--(2.9,2.9);
    \node[gray!60!black] at (2.65,2.45) {\(\overline V_+\)};
    \draw[blue!70!black,line width=1.1pt]
      plot[smooth,domain=-2.25:2.25,samples=90]
      (\x,{sqrt(\x*\x+0.95)});
    \node[blue!70!black] at (-2.25,2.65) {\(\Sigma_m^+\)};
    \draw[red!75!black,line width=2.1pt]
      plot[smooth,domain=-2.85:-2.45,samples=20]
      (\x,{sqrt(\x*\x+0.95)});
    \draw[red!75!black,line width=2.1pt]
      plot[smooth,domain=2.45:2.85,samples=20]
      (\x,{sqrt(\x*\x+0.95)});
    \fill[green!14,draw=green!45!black,line width=0.9pt]
      (1.12,1.55) ellipse (0.72 and 0.43);
    \node[green!45!black,align=center] at (1.12,0.72)
      {\(\operatorname{supp}\widehat\chi\subset N\)};
    \fill[orange!15,draw=orange!85!black,line width=0.9pt]
      (-1.12,-1.55) ellipse (0.72 and 0.43);
    \node[orange!85!black,align=center] at (-1.42,-0.72)
      {transfer of \(B^*\)};
    \fill[black] (0,0) circle (1.7pt);
    \node[below right] at (0,0) {\(\Omega\)};
  \end{tikzpicture}
  }
  \caption{Smearing a local observable with \(\chi\) restricts its
  energy-momentum transfer.  The support of \(\widehat\chi\) is chosen inside
  the isolating neighborhood of the positive mass shell, while the reflected
  support controls the adjoint operator.}
  \label{fig:haag-ruelle-spectral-transfer}
\end{figure}

\section{Creation Operators and Velocity Supports}

Let \(B\) be a regular Haag--Ruelle creator, and let
\[
  P_1B\Omega\ne0 .
\]
For a compactly supported smooth function \(h(\vec p)\) on the mass shell,
define
\[
  h_t(\vec x)=
  \int\frac{\dd^{D-1}\vec p}{(2\pi)^{D-1}}\,
  h(\vec p)\,
  \ee^{\ii\omega_{\vec p}t-\ii\vec p\cdot\vec x}
\]
and
\[
  B_t(h)=\int\dd^{D-1}\vec x\,h_t(\vec x)B(t,\vec x).
\]
The velocity support is the compact set
\[
  V(h)=\{\vec p/\omega_{\vec p}:\vec p\in\operatorname{supp}h\}.
\]
If \(V(h_i)\cap V(h_j)=\emptyset\), stationary phase and almost locality give
rapid commutator decay:
\[
  \|[B_{i,t}(h_i),B_{j,t}(h_j)]\|
  \le C_N(1+|t|)^{-N}.
\]
This inequality is the analytic expression of local propagation for separated
massive wave packets.

\begin{theorem}[Haag--Ruelle scattering for a vacuum net]
  Assume locality, Poincare covariance, the spectrum condition, a unique
  vacuum, an isolated massive one-particle shell, and enough almost-local
  operators to create a dense subspace of the one-particle sector.  For wave
  packets with pairwise disjoint velocity supports, the limits
  \[
    \lim_{t\to\pm\infty}
    B_{1,t}(h_1)\cdots B_{n,t}(h_n)\Omega
  \]
  exist in Hilbert norm.  The limits depend only on the one-particle vectors
  \(P_1B_{j,t}(h_j)\Omega\), define isometric incoming and outgoing wave
  operators from the symmetric Fock space over \(\Hilb_1\), and intertwine the
  physical Poincare representation with the free Fock representation.
\end{theorem}

The theorem is a structural bridge: local observables and spectral data
produce asymptotic particles.  It does not assume a Lagrangian, a perturbative
expansion, or a Fock-space representation of the interacting local fields.

\section{The Scattering Operator and Completeness}

Let
\[
  \Omega_{\rm in},\Omega_{\rm out}:\mathcal F_s(\Hilb_1)\to\Hilb
\]
be the wave operators supplied by the theorem.  Their ranges are closed
subspaces \(\Hilb_{\rm in}\) and \(\Hilb_{\rm out}\).  If these ranges agree
in the sector under study, the scattering operator is
\[
  S=\Omega_{\rm out}^*\Omega_{\rm in}
  \quad\text{on}\quad \mathcal F_s(\Hilb_1).
\]
It is unitary on that sector.  Asymptotic completeness is the stronger
statement that the physical sector equals the scattering range:
\[
  \Hilb_{\rm sector}=\Hilb_{\rm in}=\Hilb_{\rm out}.
\]
This is a dynamical theorem or conjecture depending on the model.  It is not
contained in locality and positivity alone.

Completeness is where genuinely nontrivial physics enters the rigorous
framework.  Bound states require additional isolated shells.  Confinement can
remove colored one-particle shells from the observable spectrum while leaving
color-singlet hadron shells.  Long-range gauge forces can replace charged mass
shells by infraparticle continua.  Massless particles require modified
asymptotic observables because velocity separation and infrared behavior are
different from the massive isolated-shell case.  These phenomena are different
spectral and asymptotic problems with their own state spaces and limiting
procedures.

\section{Relation to LSZ and Perturbation Theory}

The Haag--Ruelle construction defines the target of LSZ.  If a local field or
observable has a nonzero matrix element between the vacuum and an isolated
one-particle shell, then the two-point function has a pole or spectral atom
with residue determined by that matrix element.  LSZ is the statement that
amputating the corresponding external singularities in time-ordered
correlation functions computes matrix elements of the already-defined
operator \(S\).

Perturbation theory enters after this point.  Renormalized Feynman graphs may
compute time-ordered Green functions, and LSZ may turn those Green functions
into scattering amplitudes when the hypotheses about stable external
particles and pole residues hold.  The mathematical order is therefore
\[
\begin{gathered}
  \text{local net and spectrum}
  \longrightarrow
  \text{Haag--Ruelle wave operators}
  \longrightarrow
  S,
  \\
  S
  \longrightarrow
  \text{LSZ extraction}
  \longrightarrow
  \text{perturbative amplitudes}.
\end{gathered}
\]
This order is independent of whether the theory was discovered by a
Lagrangian, a lattice construction, an exact correlation hierarchy, or another
nonperturbative definition.
